% Part: computability
% Chapter: computability-theory
% Section: ce-sets

\documentclass[../../../include/open-logic-section]{subfiles}

\begin{document}

\olfileid{cmp}{thy}{ces}
\olsection{مجموعه‌های محاسبه‌پذیرِ برشمردنی}

\begin{defn}
مجموعه‌ای \emph{محاسبه‌پذیرِ برشمردنی} است اگر تهی باشد یا برد یک
تابع محاسبه‌پذیر باشد.
\end{defn}

\begin{history}
مجموعه‌های محاسبه‌پذیرِ برشمردنی را \emph{بازگشتیِ برشمردنی} نیز
می‌نامند. این اصطلاح اصلی است و امروزه هر دو اصطلاح، همراه با
اختصارهای «c.e.» و «r.e.»، رایج‌اند.
\end{history}

\begin{explain}
دربارهٔ معنای تعریف و دلیل تناسب اصطلاح بیندیشید. اندیشه این است که
اگر $S$ برد تابع محاسبه‌پذیر~$f$ باشد، آنگاه
\[
S = \{ f(0), f(1), f(2), \dots \},
\]
و پس می‌توان $f$ را «برشمارندهٔ» عناصر~$S$ دانست. توجه کنید که بنا بر
تعریف، لازم نیست $f$ تابعی صعودی باشد؛ یعنی لازم نیست ترتیب برشماری
صعودی باشد. حتی لازم نیست $f$ یک‌به‌یک باشد؛ یعنی تکرار در برشماری
$f(0)$، $f(1)$، $f(2)$، \dots{} از~$S$ مجاز است. برای نمونه، تابع
ثابت $f(x)=0$ مجموعهٔ $\{0\}$ را برمی‌شمارد.
\end{explain}

هر مجموعهٔ محاسبه‌پذیر، محاسبه‌پذیرِ برشمردنی است. برای دیدن این
مطلب فرض کنید $S$ محاسبه‌پذیر باشد. اگر $S$ تهی است، بنا بر تعریف
محاسبه‌پذیرِ برشمردنی است. وگرنه بگذارید $a$ عنصری دلخواه از~$S$
باشد. $f$ را چنین تعریف کنید:
\[
f(x) =
\begin{cases}
x & \text{اگر $\Char{S}(x) = 1$} \\
a & \text{در غیر این صورت.}
\end{cases}
\]
آنگاه $f$ تابعی محاسبه‌پذیر است و $S$ برد~$f$ است.

\end{document}
